\section{Characterizing Worklist-Driven Convolution}
\label{sec:char}

% ---------------------------------------------------------------------------
% 本章职责（CLAUDE.md）：架构层 insight 的分析模型与可证伪预测；不出现 WISEConv 的设计组件，每个 Insight 后用一句话指向
% 它引出的设计需求。数字：已测的直接写；待测的写成预期值并用 \tbd{} 标红，实验在旁边的注释里（"实验："开头）。
% 所有 figure 目前是占位框；caption 按目标图写。
% 措辞："follows the structure of CUTLASS's gather implicit-GEMM convolution"（用户 2026-10-31 定）；wise2 源码不链接
%   CUTLASS，不说"rewrites / derived from"。FP32 是同一结构的 SIMT kernel（CUTLASS 的 gather 卷积是 TF32）。
% ---------------------------------------------------------------------------

\providecommand{\tbd}[1]{\textcolor{red}{#1}}
\providecommand{\insight}[2]{%
  \par\smallskip\noindent\fbox{\parbox{\dimexpr\columnwidth-2\fboxsep-2\fboxrule\relax}{%
  \textbf{Insight #1.} \emph{#2}}}\par\smallskip}
\providecommand{\figbox}[2]{%
  \fbox{\parbox[c][#1][c]{\dimexpr\columnwidth-2\fboxsep-2\fboxrule\relax}{\centering\small\sffamily #2}}}

The sparsity-to-latency gap (\S\ref{sec:bg:gap}) leaves open whether a
worklist-driven convolution closes it and what it takes. Every such operator
has two stages: a construction stage that turns the mask into a worklist, and a
compute stage that runs the dense implicit GEMM with its row coordinates read
from the worklist instead of enumerated. We characterize the two stages
separately, so that each finding belongs to the execution model rather than to
a particular implementation of it. For the compute stage we ask how close its
latency comes to proportional and what that requires of the kernel
configuration (\S\ref{sec:char:affine}), why irregular rows leave the per-tile
throughput intact (\S\ref{sec:char:reduction}), and how much the order of the
worklist matters (\S\ref{sec:char:order}); for the construction stage, what it
costs relative to the computation it enables (\S\ref{sec:char:construct}).

\paragraph{Methodology.}
\textbf{Compute stage.} We study the compute kernel of \sysname
(\S\ref{sec:design:compute}), which follows the structure of CUTLASS's gather
implicit-GEMM convolution~\cite{cutlass} and adds nothing else to it: the
gather indices become worklist entries, and a grid of $\lceil n_W / B_M
\rceil$ tiles loads $B_M$ worklist rows per tile and reduces over the $k^2$
taps and $C_{in}$ channels exactly as the dense kernel does. A persistent grid
reads $n_W$ on the device, so no length reaches the host. We instantiate it
over the configuration space a dense autotuner searches: tile shape $(B_M, B_N, B_K)$, warp layout, pipeline depth, and
split-K within a CTA or across CTAs; FP16 uses \texttt{mma.sync} tensor cores
and FP32 a SIMT kernel of the same structure. Each configuration is timed in
isolation inside CUDA graphs on worklists that are given, not constructed.
\textbf{Construction stage.} We time the mask-to-worklist pass of each layer on
its own (\S\ref{sec:char:construct}). \textbf{Reference.} Dense latency
$t_{\mathrm{dense}}$ is the faster of the same kernel on the full grid and
cuDNN's non-Winograd algorithms; Winograd changes the arithmetic, and a
masked layer cannot use it. \textbf{Workloads.} The convolution layers of
YOLOv8n and YOLOv8m on WildTrack, FireFlowNet on MVSEC, and DynConv on MPII
(Table~\ref{tab:models}), driven by synthetic worklists of a given active ratio
or by the worklists their mask policies produce on real inputs (74{,}844
layer-frame samples for YOLOv8n alone). \textbf{Platforms.} RTX~3080 (68~SMs)
and Jetson AGX Orin (16~SMs) at locked clocks.

% ===========================================================================
\subsection{Proportionality and the Configuration It Requires}
\label{sec:char:affine}
% 论证链（用户 2026-10-31 定）：理想 proportional → 逐长度最优配置近似仿射 → 配置必须随工作量变（在整个 α 范围上）→
% 运行时逐长度选择的收益小，机制要么每层成本更大、要么只拿到一部分 → 按层的 activity 分布 calibrate 的单一配置拿到同一条线
% （一层实际访问的范围窄）→ 原因：两个 regime + wave 台阶 → 边界：分布漂移需要重新 calibrate。

\begin{figure}[t]
  \centering
  \figbox{3.6cm}{Fig.\ char-affine: (a) normalized latency vs.\ $\alpha$ (log), 3080 / AGX, FP32 / FP16:
  diagonal, per-length best + affine fit, dense-tuned.\\
  (b) real layer-frame samples: per-length best vs.\ distribution-calibrated, two affine fits.}
  \caption{Convolution latency normalized to dense, $t/t_{\mathrm{dense}}$,
  summed over the layers of YOLOv8n. (a)~Synthetic worklists of active ratio
  $\alpha$: the fastest configuration at each worklist length (per-length
  best) follows an affine line close to the proportional diagonal; the
  configuration a dense autotuner picks flattens at low $\alpha$.
  (b)~Worklists of real frames on four workloads: one configuration per layer,
  calibrated on separate frames of the same workload, traces the per-length
  best's affine line with an intercept larger by \tbd{0.01--0.03}.}
  \Description{Placeholder.}
  \label{fig:char-affine}
\end{figure}

\paragraph{The per-length best configuration is nearly affine.}
A worklist kernel issues only active outputs ($\eta \approx 1$, apart from the
padding of the last tile), so proportional latency, $t(\alpha) = \alpha\,
t_{\mathrm{dense}}$, is equivalent to keeping the dense throughput. We first
grant the kernel its best case and pick the fastest configuration for every
worklist length. Its latency is affine in $\alpha$ rather than proportional
(Figure~\ref{fig:char-affine}a): $t/t_{\mathrm{dense}} = 0.794\alpha + 0.258$
for FP32 and $0.643\alpha + 0.417$ for FP16 on RTX~3080, and $0.917\alpha +
0.085$ and $0.830\alpha + 0.267$ on AGX Orin, with a median per-layer $R^2$ of
0.97--0.995. Each line meets dense latency at $\alpha = 1$ (slope and intercept
sum to about one), so the whole departure from proportionality is the intercept
$b$, a latency every layer pays regardless of how many positions are active.
Table~\ref{tab:char-intercept} breaks $b$ down with hardware counters: it is
launch and synchronization latency, not memory traffic. An empty kernel spends
92\% of its stall cycles on constant-cache misses (the kernel parameters are
cold at every launch), and a single tile adds a dependent chain of barriers,
constant-cache misses, and, with split-K, memory fences, with memory-latency
stalls at only 6--27\% and the data in L2. The intercept is 31\% (FP32) and
47\% (FP16) of a median layer's dense time on RTX~3080 and 10\% and 33\% on AGX
Orin: a layer comes closer to proportional the larger it is relative to the
device.

\begin{table}[t]
  \centering\small
  \caption{Composition of the intercept $b$ on a representative layer: latency
  and dominant stall reasons of the worklist kernel with no active position, one
  tile, and the full grid (Nsight Compute, RTX~3080 and AGX Orin).}
  \label{tab:char-intercept}
  \figbox{1.6cm}{Table char-intercept: rows $n_W = 0$, one tile, $\alpha = 1$; columns latency ($\mu$s),
  SM-active time, stall shares (constant-cache miss, barrier, fixed-latency dependency, membar,
  long scoreboard), DRAM bytes; both GPUs.}
\end{table}
% 数据已有：logs/wise2_intercept/（3080：空 kernel 1.0–1.2 µs、SM 活跃约 0.5 µs、92% imc_miss；单 tile 再加 2.6–3.3 µs
%   FP32 / 1.4–2.2 µs FP16；AGX 空 kernel 2.2–2.6 µs，单 tile 再加 FP32 8–19 / FP16 3.5–6.5 µs）。

\begin{figure}[t]
  \centering
  \figbox{3.0cm}{Fig.\ char-persm: slowdown of a fixed configuration over the per-length best vs.\ work per SM
  (log), dense-tuned and activity-agnostic single configuration, 3080 and AGX overlaid.}
  \caption{Slowdown of a fixed configuration over the per-length best across
  the synthetic sweep of Figure~\ref{fig:char-affine}a, binned by the work per SM
  of the fixed configuration (its tiles times its K splits, divided by the SM
  count). Both GPUs fall on one curve: a configuration is near-optimal only in
  the regime of work per SM it suits.}
  \Description{Placeholder.}
  \label{fig:char-persm}
\end{figure}

\paragraph{Over the full range of activity, the best configuration changes.}
No single configuration stays near this line from $\alpha \to 0$ to $\alpha =
1$. The configuration a dense autotuner picks, the fastest at $\alpha = 1$,
runs 2.33$\times$ (FP32) and 1.64$\times$ (FP16) slower than the per-length best
when it gives an SM less than a quarter of a tile on RTX~3080, and 4.00$\times$
and 1.92$\times$ on AGX Orin, and matches the best above four tiles per SM
(Figure~\ref{fig:char-persm}). Large tiles amortize each weight load over many
rows; when rows are few, smaller tiles and split-K expose enough independent
work to occupy the SMs and shorten the chain each one executes. Picking the
configuration with the lowest expected latency over $\alpha$ uniform on $(0,
1]$ only moves the loss: it is still 1.61$\times$ and 1.39$\times$ slower below a
quarter of a tile per SM on RTX~3080, and 11.6--13.5\% slower at $\alpha = 1$.
A layer therefore needs a configuration that depends on its activity. Whether
it needs one per worklist length depends on how much of this range a layer
actually visits, which we examine next.

\begin{figure}[t]
  \centering
  \figbox{3.0cm}{Fig.\ char-mechanism: CDF over layers of the expected per-frame gain of per-length selection
  over the layer's best single configuration ($\mu$s); vertical lines: per-layer cost of host readback,
  updatable node, conditional node; markers: realized gain of in-kernel dispatch and of lagged graph selection.}
  \caption{Per-layer gain that selecting the configuration for every worklist
  length could deliver over the best single configuration of the layer (CDF over
  the layers of four workloads, RTX~3080 and AGX Orin), against what each runtime
  mechanism costs or realizes: host readback and launch, a device-updatable
  kernel node, and a conditional graph node pay a fixed cost per layer; in-kernel
  dispatch among fused configurations and a CUDA graph chosen per frame from
  lagged activity pay almost nothing but realize only part of the gain.}
  \Description{Placeholder.}
  \label{fig:char-mechanism}
\end{figure}

\paragraph{Selecting per length gains little, and every mechanism pays for it.}
Against the best single configuration for the layer's own frames, choosing per
worklist length saves a median layer only \tbd{0.6}\,$\mu$s per frame on
RTX~3080 and \tbd{1.5}\,$\mu$s on AGX Orin (Figure~\ref{fig:char-mechanism}).
Capturing this gain at run time is costly, because the length that drives the
choice exists only on the device and changes at every layer of every frame.
Mechanisms that decide per layer pay a fixed cost per layer: reading the length
back serializes each layer behind a host synchronization, a device-updatable
kernel node rewrites the next launch at \tbd{0.9}\,$\mu$s on RTX~3080 and
\tbd{2.2--3.3}\,$\mu$s on AGX Orin, and a conditional graph node selects a
subgraph at \tbd{17--19}\,$\mu$s on AGX Orin; only \tbd{X\%} of the layers gain
more than the cheapest of them costs. Mechanisms that avoid the per-layer cost
give up part of the gain. Fusing the candidate configurations into one kernel
and branching inside it decides for free, but the fused kernel carries the
register and code footprint of all its paths and realizes \tbd{30--60\%} of the
gain. Choosing one of a few pre-captured CUDA graphs per frame from the activity
of past frames decides on the host without synchronization, but a frame-level
prediction cannot follow layers whose activity varies independently, and it
realizes \tbd{50\%} of the gain on YOLOv8n and \tbd{10\%} on DynConv.
% 实验 E2（需 GPU，约半小时 / 设备）：
%   收益：每层 per-frame oracle 相对该层在自身帧上最优的单一配置的期望增益（µs），隔离 sweep × 真实 trace，四个 workload。
%     0.6 µs 是粗估（3080 FP32 YOLOv8n，in-situ：(S − O) / 层数 ≈ 40.6 / 63）。
%   成本：当前 binary 重测：host 读回 + 分支 launch、device-updatable node（m0_probe）、conditional SWITCH node。
%     0.9 / 2.2–3.3 / 17–19 µs 是第一代原型的数，必须重测后才能保留。
%   实现比例：in-kernel dispatch 用 in-situ record（3080：D 只拿到 O 的 30–60%，static_vs_dynamic.py）；lagged graph 用
%     bucket_static.py（3080 YOLOv8n K = 4：−2.8% vs 每帧 oracle −5.1%，约 55%；DynConv 帧级理想 −0.6% vs −1.5%）。

\paragraph{One calibrated configuration per layer reaches the same line.}
The range a layer visits is far narrower than $(0, 1]$, so a configuration
suited to that range suffices. We calibrate one configuration per layer to its
activity distribution, taking the configuration with the lowest mean latency on
calibration frames, and evaluate it on separate frames of the same workload (a
held-out camera for WildTrack, held-out clips otherwise). On real worklists its
affine fit shares the per-length best's slope, and its intercept is larger by
only \tbd{0.01--0.03} of dense time on both GPUs and precisions
(Figure~\ref{fig:char-affine}b); summed over a model it is within
\tbd{2--4\%} of the per-length best on YOLOv8n, YOLOv8m, and FireFlowNet and
within \tbd{1\%} on DynConv. DynConv is the hardest case for a single choice:
its learned gates vary a layer's active ratio by up to 9$\times$ between the
tenth and ninetieth percentile frame, nearly independently of the other
layers. The configuration a dense autotuner picks, in contrast, is 42--94\%
slower than the calibrated one on the YOLO models: what matters is calibrating
to the activity a layer sees, not adapting to each frame of it.
% 实验 E1（离线，数据已有）：隔离 sweep（standalone persistent tile，与 (a) 同一配置空间）× 真实 trace，每个
%   (层, 帧) 的 t/t_dense；per-length best、calibrated（交叉验证）、dense-tuned；两条仿射拟合 (a, b, R²)、Δb、模型级
%   Σt 比；四个 workload × 3080 / AGX × FP32 / FP16（AGX YOLOv8m 待部署）。
% 风险：隔离 sweep 上 per-frame oracle 相对 calibrated 在 YOLO 上是 −9 至 −11%（logs/wise2_static/，s = 1）；in-situ 只在
%   4 个 static 候选里比是 −1.3 至 −4.5%。若 E1 的 Δb 明显大于 0.03，本段改为"差距存在，但小于任何机制能拿到的部分"，
%   论证重心落到上一段。42–94%：隔离 sweep，anchor vs static（YOLOv8n FP32 / FP16 +66% / +42%，YOLOv8m +94% / +79%）。
%   DynConv 9×：单层 α 的 P10–P90 例如 0.02–0.18（trace_fp32，3080）；跨层相关中位数 0.12–0.42。

\begin{figure}[t]
  \centering
  \figbox{3.0cm}{Fig.\ char-wave: (a) $t(n)$ of the calibrated configuration and the per-length best on one
  sub-wave and one multi-wave layer, with the layer's $n_W$ distribution; (b) regret vs.\ wave phase;
  (c) per-layer P10--P90 span of the work per SM.}
  \caption{Why one configuration suffices. (a)~Latency against worklist length
  for a layer whose worklists fill less than one wave and for a layer that spans
  several, with the layer's distribution of lengths shaded: below one wave the
  latency is nearly flat, above it a staircase. (b)~Regret of the calibrated
  configuration over the per-length best against the calibrated configuration's
  wave phase: in the multi-wave regime it concentrates just past a wave
  boundary. (c)~Span of each layer's work per SM between its tenth and ninetieth
  percentile frame.}
  \Description{Placeholder.}
  \label{fig:char-wave}
\end{figure}

\paragraph{Why one configuration suffices.}
A configuration's latency has two regimes in the worklist length
(Figure~\ref{fig:char-wave}a). While its tiles fit in one wave of co-resident
CTAs, every CTA runs at most one tile and the latency is the single-tile chain
plus the launch, nearly flat in $n_W$; the best configuration here is the one
with the shortest chain at the layer's shape, which does not move with $n_W$.
Beyond one wave the latency is a staircase, $t_q(n) \approx
\lceil \mathrm{tiles}_q(n) / W_q \rceil\, t_{\mathrm{tile},q} + t_{0,q}$ for
$W_q$ co-resident CTAs, and configurations with different tile shapes place
their steps at different lengths: the per-length best beats a fixed
configuration mainly just past one of its steps, where a partial wave costs a
full tile time (Figure~\ref{fig:char-wave}b). These steps account for
\tbd{70\%} of the per-length gain; the rest comes from layers whose lengths
straddle the two regimes. Most layers do not: between its tenth and ninetieth
percentile frame, a median layer's work per SM spans a factor of \tbd{4}
(Figure~\ref{fig:char-wave}c), and \tbd{X\%} of the layers stay on one side of a
wave throughout. A layer thus spends nearly all its frames in one regime, where
one configuration is best up to the steps.
% 实验 E3（离线，数据已有）：(a) 两个代表层的 t(n) 曲线与 n_W 直方图；(b) 每层 regret(n) 对 calibrated 配置的 wave 相位
%   （tiles·splits / (SM · occ) mod 1；occ 用 occupancy API 或由台阶位置推出）；(c) 阶梯残差分解：每个配置的 t(n) 拆成
%   该层 n 范围内的仿射趋势 + 阶梯残差，oracle 增益中来自残差的比例（文中 70%），以及跨 regime 的层的贡献；(d) 每层 n_W 的
%   P10–P90 跨度（tiles/SM，文中 4×）与始终在一个 wave 内 / 外的层比例。
% 注意：YOLOv8n 95%（3080）/ 73%（AGX）的样本 < 1 tile/SM，主要在 sub-wave regime；需确认 sub-wave 内"近似平"成立
%   （CLAUDE.md：同一 wave 内时间也随活跃 CTA 增长），否则改写 sub-wave 的机制描述。

\paragraph{The distribution, not the frame, sets the configuration.}
Calibration does depend on the distribution it sees. Scaling every layer's
worklist lengths by a factor between $1/4$ and $4$ makes the configuration
calibrated at the original scale 3--24\% slower than one re-calibrated at the
new scale, which in turn stays within \tbd{1--3\%} of the per-length best. Real
shifts behave the same way: a configuration calibrated on one MVSEC sequence
loses \tbd{X\%} on a sequence with \tbd{4$\times$} its activity. A deployment
must therefore calibrate to the activity it will see and recalibrate when that
activity changes substantially, but within a distribution it needs no per-frame
adaptation.
% 实验 E4：合成漂移已有（static_vs_dynamic.py 的 shift，3080 四个 workload：沿用 3–24%；重选距 oracle：YOLO 3.5–12%，
%   DynConv 0.3–2%，需在 E1 口径下复核文中的 1–3%）。真实漂移：held-out split 上跑 insitu_calib.py trace（FireFlowNet
%   outdoor_day1 vs calib 的 indoor_flying1；WildTrack randpos vs calib），calib 上选的配置在 held-out 上的 regret 与重选对比。

\insight{1}{With the configuration chosen for every worklist length, a
worklist-driven implicit GEMM is nearly affine in the active ratio, but that
choice gains little over one configuration calibrated to a layer's activity
distribution, and every runtime mechanism either costs more per layer or
realizes only part of it. The calibrated configuration reaches the same line
with a slightly larger intercept, because a layer's work per SM stays in one
regime across frames, where per-length selection only exploits
wave-quantization steps.}

The compute stage (\S\ref{sec:design:compute}) therefore provides a
configuration space rich in fine decompositions, calibrates one configuration
per layer to the deployment's activity, and has the kernel consume the
device-resident worklist length directly.

% ===========================================================================
\subsection{Irregularity Stays Outside the Reduction Loop}
\label{sec:char:reduction}

\begin{figure}[t]
  \centering
  \figbox{3.0cm}{Fig.\ char-reduction: per representative layer, (a) global-load sectors per request,
  worklist vs.\ dense; (b) DRAM bytes: measured, active-halo-union prediction, $k^2 n_W C_{in}$.}
  \caption{Per-tile memory behavior of the worklist kernel on seven
  representative YOLOv8n layers (RTX~3080 shown; AGX Orin is the same).
  (a)~Sectors per global-load request stay close to the dense kernel's, except
  for the stem, whose 32-byte rows cannot fill a request. (b)~DRAM traffic
  equals the union of the active positions' input halos plus the weights, not
  the $k^2 n_W C_{in}$ a gather-scatter path moves.}
  \Description{Placeholder.}
  \label{fig:char-reduction}
\end{figure}

The affine line meets dense latency at $\alpha = 1$, so an active position
costs the worklist kernel what it costs the dense kernel. The reason lies in
where the worklist enters the kernel. An implicit
GEMM draws on three levels of reuse: contiguous channels within a row (memory
coalescing), a weight tile shared by the $B_M$ rows of a CTA (arithmetic
intensity), and overlapping receptive fields of neighboring outputs (L1 and L2
hits). The worklist changes only the base address of each row. A tile issues
$B_M$ index loads in its prologue, ahead of the multistage pipeline, and the
reduction loop over $k^2$ taps and channel slices keeps the dense kernel's
affine addressing. The first two levels of reuse are therefore untouched.
Counters confirm it: a worklist tile executes the same instructions as a dense
tile (within 0.5\%, with identical FMA, tensor, and global-load counts), and
global loads keep 12.8--14.3 sectors per request against 14.5--15.5 in the
dense kernel for FP32 (8.6--12.7 against 11.0--15.4 for FP16;
Figure~\ref{fig:char-reduction}a). The one exception is the predicted boundary:
the stem layer has eight padded input channels, a 32-byte row that fills only
3.7 sectors per request in both kernels. DRAM traffic is set by the union of
the active halos: measured bytes are 1.00--1.11$\times$ that prediction, while
gathering each position's receptive field separately would move 2--7$\times$
more on 3$\times$3 layers (Figure~\ref{fig:char-reduction}b). Point-cloud sparse
convolution differs exactly here: its input is sparse too, so every reduction
step depends on a rulebook lookup.

% ===========================================================================
\subsection{Worklist Order Barely Matters}
\label{sec:char:order}

\begin{figure}[t]
  \centering
  \figbox{3.0cm}{Fig.\ char-order: distribution over layer-frame samples of latency relative to the
  construction order, for sorted / shuffled / uniform worklists, 3080 and AGX; inset: DRAM, L1-hit, L1-to-L2
  traffic on representative layers.}
  \caption{Latency of four worklists with the same length and tile on real
  frames, relative to the order construction emits: sorted in global dense
  order, shuffled, and uniform (the same count of positions spread uniformly,
  removing the mask's spatial clustering). Inset: counters on representative
  layers; order changes L1 hits, never DRAM bytes.}
  \Description{Placeholder.}
  \label{fig:char-order}
\end{figure}

The third level of reuse, overlapping receptive fields, depends on which
positions share a CTA and therefore on the order of the worklist. It matters
little (Figure~\ref{fig:char-order}). Across 3{,}560 layer-frame samples, a
globally sorted worklist runs at 0.999$\times$ the latency of the order
construction emits (median) and a shuffled one at 1.004--1.007$\times$; summed
over the model, every variant is within 2.2\% on RTX~3080. Counters show why:
the construction order, sorted, and shuffled worklists move identical DRAM bytes, and order only
shifts reuse between L1 (6--46\% hits for the construction order, 3--26\%
shuffled) and L2, which absorbs the rest. Order becomes visible only where L2
bandwidth is scarce relative to the work: on AGX Orin, shuffling slows the
early layers with many positions and little computation each by
1.10--1.33$\times$, because the extra L1-to-L2 traffic (7--70\% more) surfaces on
a GPU with little L2 bandwidth. Even there, a tile-local order recovers all of
the global order's benefit (sorted against construction order:
0.993--0.999$\times$).

\insight{2}{Worklist order never changes DRAM traffic, which the union of the
active halos sets; it only shifts reuse between L1 and L2. A random order costs
about 1\% except on low-intensity layers of a GPU with little L2 bandwidth, and
a tile-local order already recovers everything a global order would.}

Construction therefore needs to preserve only tile-local order
(\S\ref{sec:design:construct}).

% ===========================================================================
\subsection{Construction Cost Does Not Shrink with Activity}
\label{sec:char:construct}

\begin{figure}[t]
  \centering
  \figbox{3.0cm}{Fig.\ char-construct: per-frame construction share of convolution time vs.\ frame
  $\alpha$, 3080 / AGX, FP32 / FP16; inset: per-layer construction vs.\ compute by $C_{in}$.}
  \caption{Share of convolution time spent building worklists, per YOLOv8n
  frame, against the frame's active ratio. Construction costs a near-constant
  time per layer, so its share rises as activity falls and as tensor cores
  shorten the compute.}
  \Description{Placeholder.}
  \label{fig:char-construct}
\end{figure}

A worklist must exist before the computation it enables, and building it scans
the mask grid of the layer's geometry, whose size does not depend on activity.
Its cost is therefore nearly constant per layer: 1.7--3.2\,$\mu$s on RTX~3080
and 4.2--23.6\,$\mu$s on AGX Orin, dominated by a fixed launch cost (about
1.8\,$\mu$s and 4\,$\mu$s) and visible mask scanning only on the 640$\times$640
stem. The computation it enables shrinks with activity and channel count.
Over a YOLOv8n frame, 52 constructions (the other 11 layers reuse their
input's worklist) take 93\,$\mu$s on RTX~3080 against
397\,$\mu$s (FP32) and 232\,$\mu$s (FP16) of computation, a median share of
19\% and 29\%; on AGX Orin they take 258\,$\mu$s, 11\% and 27\%. The share
tracks the frame's activity (correlation $-0.94$ and $-0.96$ on RTX~3080,
$-0.81$ and $-0.93$ on AGX Orin): it reaches 24\% and 33\% in the
least active quarter of frames on RTX~3080 and falls to 12\% and 23\% in the
most active quarter (Figure~\ref{fig:char-construct}). Thin layers pay the most,
because the computation scales with $C_{in}$ and the construction does not
(construction reaches \tbd{0.48} of compute below 32 input channels against
\tbd{0.36} above 96).
% 0.48 / 0.36 是 construction 改进（2026-10-31）之前的 3080 FP32 数，需用当前实现重算（construction_cost.py 的逐层输出）。

\insight{3}{Building a worklist costs a near-constant time per layer,
independent of activity and channel count, so its share of the convolution
grows exactly where sparsity is most useful.}

Construction must therefore touch only the mask and be shared wherever layers
can use the same worklist (\S\ref{sec:design:construct},
\S\ref{sec:design:reuse}).
